\documentclass[12pt,a4paper]{amsart}

\usepackage[ngerman]{babel}
\usepackage{amsmath,amssymb,amsthm}
\usepackage{mathtools}
\usepackage[colorlinks=true,linkcolor=blue,citecolor=blue]{hyperref}
\usepackage{geometry}
\geometry{margin=2.5cm}

% -------------------------------------------------------
% Theorem-Umgebungen
% -------------------------------------------------------
\newtheorem{theorem}{Satz}[section]
\newtheorem{lemma}[theorem]{Lemma}
\newtheorem{korollar}[theorem]{Korollar}
\newtheorem{proposition}[theorem]{Proposition}
\theoremstyle{definition}
\newtheorem{definition}[theorem]{Definition}
\newtheorem{bemerkung}[theorem]{Bemerkung}

% -------------------------------------------------------
% Eigene Makros
% -------------------------------------------------------
\newcommand{\N}{\mathbb{N}}
\newcommand{\Z}{\mathbb{Z}}
\newcommand{\Q}{\mathbb{Q}}
\newcommand{\Fp}{\mathbb{F}_p}
\DeclareMathOperator{\ord}{ord}
\DeclareMathOperator{\ggT}{ggT}

% -------------------------------------------------------
\begin{document}
% -------------------------------------------------------

\title{Der Satz von Wilson und seine Verallgemeinerungen:\\
       Ein elementarer und vollständiger Beweis}

\author{Michael Fuhrmann}
\address{Specialist-Mathematics-Projekt, Build 62}
\email{specialist-maths@localhost}
\date{\today}

\begin{abstract}
Der Satz von Wilson besagt, dass eine ganze Zahl $p \geq 2$ genau dann eine
Primzahl ist, wenn $(p-1)! \equiv -1 \pmod{p}$ gilt.
Wir präsentieren einen vollständigen, elementaren Beweis dieses klassischen
Ergebnisses mithilfe grundlegender Eigenschaften modularer Arithmetik und
der Struktur der multiplikativen Gruppe $(\Z/p\Z)^*$.

Darüber hinaus beweisen wir vier natürliche Verallgemeinerungen:
\begin{enumerate}
  \item[(i)]   Den exakten Rest von $(n-1)!$ modulo $n$ für zusammengesetzte $n$,
  \item[(ii)]  einen Wilson-artigen Satz für Primzahlpotenzen $p^k$,
  \item[(iii)] den Wilson-Quotienten $W_p = \bigl((p-1)!+1\bigr)/p$
               und seinen Bezug zu Wilson-Primzahlen,
  \item[(iv)]  eine Produktformel über alle Elemente einer endlichen abelschen
               Gruppe (Allgemeiner Wilson-Satz).
\end{enumerate}
Alle Beweise sind elementar und benötigen keine Hilfsmittel jenseits der
Zahlentheorie des Grundstudiums.
\end{abstract}

\maketitle
\tableofcontents

% ================================================================
\section{Einleitung}
% ================================================================

Das heute als Satz von Wilson bekannte Ergebnis erscheint erstmals in einem
Brief von Edward Waring aus dem Jahr 1770, der es John Wilson zuschrieb.
Der erste veröffentlichte Beweis stammte von Lagrange (1771) mithilfe
von Polynomkongruenzen über $\Fp$.

\begin{theorem}[Wilson 1770; Lagrange 1771]\label{satz:wilson}
  Eine ganze Zahl $p \geq 2$ ist genau dann eine Primzahl, wenn
  \[
    (p-1)! \;\equiv\; -1 \pmod{p}.
  \]
\end{theorem}

Der Satz liefert eine vollständige Charakterisierung der Primzahlen,
ist aber als Primalitätstest unpraktisch, da die Berechnung von $(p-1)!$
exponentiell viele Multiplikationen erfordert.
Seine Bedeutung liegt in der strukturellen Einsicht über die
multiplikative Gruppe modulo einer Primzahl.

Im gesamten Artikel stehen alle Kongruenzen modulo der nach $\pmod{}$
geschriebenen Zahl. Wir schreiben $a \mid b$ für Teilbarkeit und~$p$
ausschließlich für eine Primzahl.

% ================================================================
\section{Beweis des Satzes von Wilson}
% ================================================================

\subsection{Die „Wenn"-Richtung: Primzahlen erfüllen $(p-1)! \equiv -1$}

Der Beweis stützt sich auf eine Schlüsselbeobachtung über selbst-inverse
Elemente.

\begin{lemma}[Selbst-inverse Elemente modulo $p$]\label{lem:selbstinvers}
  Sei $p$ eine Primzahl und $a \in \{1, 2, \ldots, p-1\}$.
  Dann gilt $a^2 \equiv 1 \pmod{p}$ genau dann, wenn $a \in \{1, p-1\}$.
\end{lemma}

\begin{proof}
  Aus $a^2 \equiv 1 \pmod{p}$ folgt $p \mid a^2 - 1 = (a-1)(a+1)$.
  Da $p$ eine Primzahl ist, gilt $p \mid a-1$ oder $p \mid a+1$,
  also $a \equiv 1$ oder $a \equiv -1 \pmod{p}$.
  Beide Werte liegen in $\{1, p-1\}$ (für $p \geq 3$;
  für $p = 2$ gilt $1 \equiv -1$).
\end{proof}

\begin{proof}[Beweis von Satz~\ref{satz:wilson}, „Wenn"-Richtung]
  Sei $p$ eine Primzahl.  Für $p = 2$: $(2-1)! = 1 \equiv -1 \pmod{2}$. \checkmark

  Sei nun $p \geq 3$.
  Jedes Element $a \in \{1, \ldots, p-1\}$ besitzt ein eindeutiges
  multiplikatives Inverses $a^{-1} \in \{1, \ldots, p-1\}$ modulo~$p$
  (da $(\Z/p\Z)^*$ eine Gruppe der Ordnung $p-1$ ist).

  Gemäß Lemma~\ref{lem:selbstinvers} sind $1$ und $p-1$ ihre eigenen Inversen.
  Alle anderen Elemente $a \in \{2, \ldots, p-2\}$ erfüllen $a^{-1} \neq a$,
  bilden also $(p-3)/2$ disjunkte Paare $\{a, a^{-1}\}$
  (beachte: $p-3$ ist gerade, da $p$ ungerade).

  Wir multiplizieren alle Elemente von $\{1, \ldots, p-1\}$:
  \[
    (p-1)!
    = 1 \cdot \Bigl[\prod_{\substack{a=2 \\ a \neq a^{-1}}}^{p-2} a\Bigr]
      \cdot (p-1)
    = 1 \cdot \Bigl[\prod_{\{a,a^{-1}\}} a \cdot a^{-1}\Bigr] \cdot (p-1).
  \]
  Jedes Paar $\{a, a^{-1}\}$ liefert den Faktor $a \cdot a^{-1} \equiv 1 \pmod{p}$.
  Daher gilt
  \[
    (p-1)! \;\equiv\; 1 \cdot 1^{(p-3)/2} \cdot (p-1)
             \;\equiv\; p-1 \;\equiv\; -1 \pmod{p}. \qedhere
  \]
\end{proof}

\subsection{Die „Nur wenn"-Richtung: zusammengesetzte Zahlen versagen}

\begin{proof}[Beweis von Satz~\ref{satz:wilson}, „Nur wenn"-Richtung]
  Sei $n \geq 2$ zusammengesetzt.  Wir zeigen $(n-1)! \not\equiv -1 \pmod{n}$.

  Schreibe $n = ab$ mit $1 < a \leq b < n$.

  \textbf{Fall 1: $a < b$.}
  Dann erscheinen $a$ und $b$ als verschiedene Faktoren in
  $(n-1)! = 1 \cdot 2 \cdots (n-1)$,
  also gilt $n = ab \mid (n-1)!$, d.\,h.\ $(n-1)! \equiv 0 \pmod{n} \neq -1$.

  \textbf{Fall 2: $a = b$, d.\,h.\ $n = a^2$ mit $a \geq 2$.}
  \begin{itemize}
    \item Falls $a \geq 3$: Dann sind $a$ und $2a$ beide in $\{1, \ldots, n-1\}$
      (da $2a = 2\sqrt{n} < n$ für $n \geq 9$), und $a \cdot 2a = 2n$,
      also $n \mid (n-1)!$; folglich $(n-1)! \equiv 0 \pmod{n}$.
    \item Falls $a = 2$: $n = 4$, und $(4-1)! = 6 \equiv 2 \pmod{4} \neq -1$.
  \end{itemize}
  In allen Fällen gilt $(n-1)! \not\equiv -1 \pmod{n}$. \qedhere
\end{proof}

% ================================================================
\section{Der Rest von $(n-1)!$ für zusammengesetzte $n$}
% ================================================================

\begin{theorem}\label{satz:zusammengesetzt}
  Sei $n \geq 6$ zusammengesetzt.  Dann gilt $(n-1)! \equiv 0 \pmod{n}$.
\end{theorem}

\begin{proof}
  Sei $n \geq 6$ zusammengesetzt.  Wir unterscheiden zwei Fälle.

  \textbf{Fall 1: $n$ ist kein Quadrat.}
  Dann besitzt $n$ eine Zerlegung $n = ab$ mit $1 < a < b < n$.
  Beide Faktoren $a$ und $b$ erscheinen als verschiedene Faktoren in $(n-1)!$,
  also gilt $n = ab \mid (n-1)!$, d.\,h.\ $(n-1)! \equiv 0 \pmod{n}$.

  \textbf{Fall 2: $n = a^2$ ist eine Quadratzahl.}
  Da $n \geq 6$ und $n = a^2$, gilt $a \geq 3$
  (der Fall $a = 2$, $n = 4$ ist ausgeschlossen).
  Dann sind $a$ und $2a$ beide in $\{1, \ldots, n-1\}$
  (da $2a < a^2 = n$ für $a \geq 3$) und $a \neq 2a$.
  Ihr Produkt $a \cdot 2a = 2a^2 = 2n$ erscheint in $(n-1)!$,
  also gilt $n \mid (n-1)!$, d.\,h.\ $(n-1)! \equiv 0 \pmod{n}$.
\end{proof}

\begin{korollar}
  Für $n \geq 2$ nimmt der Rest $(n-1)! \bmod n$ genau dann den Wert $-1$
  (d.\,h.\ $n-1$) an, wenn $n$ eine Primzahl ist.
\end{korollar}

% ================================================================
\section{Wilson-Satz für Primzahlpotenzen}
% ================================================================

Das folgende Lemma ist das Kernwerkzeug; wir formulieren es vorab.

\begin{lemma}[Produkt aller Elemente einer endlichen abelschen Gruppe]%
  \label{lem:produkt_abelsch}
  Sei $G$ eine endliche abelsche Gruppe und $S = \{g \in G : g^2 = e\}$.
  \begin{enumerate}
    \item[(a)] $|S| = 1$ (nur $e$): $\prod_{a \in G} a = e$.
    \item[(b)] $|S| = 2$ ($S = \{e, \tau\}$ mit eindeutigem $\tau$):
      $\prod_{a \in G} a = \tau$.
    \item[(c)] $|S| \geq 4$: $\prod_{a \in G} a = e$.
  \end{enumerate}
\end{lemma}

\begin{proof}
  Wir ordnen jedem $a \in G$ sein Inverses $a^{-1}$ zu.  Gilt $a \neq a^{-1}$,
  liefert das Paar $a \cdot a^{-1} = e$.  Die ungepaarten Elemente sind genau
  die Elemente von $S$.  Also ist $\prod_{a \in G} a = \prod_{s \in S} s$.
  In Fall (a): $e$.  In Fall (b): $e\cdot\tau = \tau$.
  In Fall (c): $S \cong (\Z/2\Z)^k$ für $k \geq 2$, und in jeder Koordinate
  tragen genau $2^{k-1}$ (gerade viele) Elemente eine $1$ bei, also ist
  die Koordinatensumme $0$, d.\,h.\ das Produkt ist $e$.
\end{proof}

\begin{theorem}[Wilson-Satz für $p^k$]\label{satz:wilson_primzahlpotenz}
  Sei $p$ eine ungerade Primzahl und $k \geq 1$.  Dann gilt
  \[
    \prod_{\substack{j=1 \\ \ggT(j,p^k)=1}}^{p^k} j
    \;\equiv\; -1 \pmod{p^k}.
  \]
  Für $p = 2$: Das Produkt ist $-1 \pmod{2}$ ($k=1$), $-1 \pmod{4}$ ($k=2$)
  und $+1 \pmod{2^k}$ für alle $k \geq 3$.
\end{theorem}

\begin{proof}
  \textbf{Ungerades $p$, beliebiges $k \geq 1$.}
  $G = (\Z/p^k\Z)^*$ ist zyklisch von Ordnung $\varphi(p^k) = p^{k-1}(p-1)$
  (Primitivwurzeln existieren modulo $p^k$; vgl.\ \cite{Ireland1990}, Kap.~4).
  Die selbst-inversen Elemente erfüllen $x^2 \equiv 1 \pmod{p^k}$, also
  $p^k \mid (x-1)(x+1)$.  Da $p$ ungerade ist, teilt $p$ nicht beide Faktoren
  (Differenz $= 2$, teilerfremd zu $p$), also gilt $p^k \mid x-1$ oder $p^k \mid x+1$,
  d.\,h.\ $x \equiv \pm 1 \pmod{p^k}$.  Damit $|S| = 2$ mit $\tau = -1$.
  Nach Lemma~\ref{lem:produkt_abelsch}(b): Produkt $= -1 \pmod{p^k}$.

  \textbf{$p = 2$, $k = 1$.}  Einziges Element: $1$.  Produkt $= 1 \equiv -1 \pmod{2}$.

  \textbf{$p = 2$, $k = 2$.}  Einheiten: $\{1,3\}$.  Produkt $= 3 \equiv -1 \pmod{4}$.

  \textbf{$p = 2$, $k \geq 3$.}
  Die Lösungen von $x^2 \equiv 1 \pmod{2^k}$ sind $x \equiv \pm 1 \pmod{2^{k-1}}$,
  also vier Lösungen modulo $2^k$:
  $S = \{1,\;2^{k-1}-1,\;2^{k-1}+1,\;2^k-1\}$, $|S| = 4$.
  Direkt: $\prod_{s \in S} s = 1\cdot(2^{k-1}-1)(2^{k-1}+1)\cdot(-1) = -(2^{2k-2}-1)$.
  Da $k \geq 3$ impliziert $2k-2 \geq k$, gilt $2^{2k-2} \equiv 0 \pmod{2^k}$,
  also $\prod_{s\in S}s \equiv -(-1) = 1 \pmod{2^k}$.
  Nach Lemma~\ref{lem:produkt_abelsch}(c): Produkt $= 1 \pmod{2^k}$.
\end{proof}

% ================================================================
\section{Der Wilson-Quotient und Wilson-Primzahlen}
% ================================================================

\begin{definition}\label{def:wilson_quotient}
  Für eine Primzahl $p$ ist der \emph{Wilson-Quotient} definiert als
  \[
    W_p \;=\; \frac{(p-1)!\,+\,1}{p} \;\in\; \Z.
  \]
  (Dies ist eine ganze Zahl nach dem Satz von Wilson.)
  Eine Primzahl $p$ heißt \emph{Wilson-Primzahl}, wenn $p^2 \mid (p-1)! + 1$,
  äquivalent $p \mid W_p$.
\end{definition}

\begin{proposition}\label{prop:wilson_quotient}
  Der Wilson-Quotient erfüllt:
  \[
    W_p \;\equiv\; \frac{(p-1)!+1}{p} \pmod{p}.
  \]
  Insbesondere ist $p$ genau dann eine Wilson-Primzahl, wenn $W_p \equiv 0 \pmod{p}$.
\end{proposition}

\begin{proof}
  Dies folgt unmittelbar aus der Definition: $W_p \in \Z$ und die Bedingung
  $p^2 \mid (p-1)!+1$ ist äquivalent zu $p \mid W_p$.
\end{proof}

\begin{bemerkung}
  Die einzigen bekannten Wilson-Primzahlen sind $5$, $13$ und $563$.
  Es wird vermutet (aber nicht bewiesen), dass unendlich viele Wilson-Primzahlen
  existieren.  Für diese drei Primzahlen gilt:
  \[
    (5-1)! + 1 = 25 = 5^2, \quad
    W_{13} \equiv 0 \pmod{13}, \quad
    W_{563} \equiv 0 \pmod{563}.
  \]
\end{bemerkung}

% ================================================================
\section{Allgemeiner Wilson-Satz für endliche abelsche Gruppen}\label{satz:allgemein_wilson}
% ================================================================

\begin{theorem}[Allgemeiner Wilson-Satz]\label{satz:allgemein_wilson}
  Sei $(G, \cdot)$ eine endliche abelsche Gruppe.  Dann gilt
  \[
    \prod_{g \in G} g
    \;=\;
    \begin{cases}
      e        & \text{falls } G \text{ kein Element der Ordnung } 2 \text{ hat}, \\
      \tau      & \text{falls } G \text{ genau ein Element } \tau \text{ der Ordnung } 2 \text{ hat}, \\
      e        & \text{falls } G \text{ mehr als ein Element der Ordnung } 2 \text{ hat}.
    \end{cases}
  \]
  Für $G = (\Z/p\Z)^*$ mit einer ungeraden Primzahl $p$ ist das eindeutige
  Element der Ordnung $2$ gleich $-1$, was $(p-1)! \equiv -1 \pmod{p}$ liefert.
\end{theorem}

\begin{proof}
  Wir ordnen jedem $g \in G$ sein Inverses $g^{-1}$ zu.
  Falls $g \neq g^{-1}$, liefert das Paar $\{g, g^{-1}\}$ den Beitrag
  $g \cdot g^{-1} = e$ zum Produkt.
  Nur die selbst-inversen Elemente (Ordnung $1$ oder $2$) bleiben übrig.

  Die Menge $S = \{g \in G : g^2 = e\}$ ist eine Untergruppe von $G$.
  \begin{itemize}
    \item $|S| = 1$: $S = \{e\}$, das Produkt ist $e$.
    \item $|S| = 2$: $S = \{e, \tau\}$, das Produkt ist $e \cdot \tau = \tau$.
    \item $|S| \geq 4$: $S \cong (\Z/2\Z)^k$ für $k \geq 2$.
      Das Produkt aller Elemente von $(\Z/2\Z)^k$ ist $0$ (im additiven Sinne),
      also $e$ in multiplikativer Notation.
  \end{itemize}
  Für $G = (\Z/p\Z)^*$ ($p$ ungerade Primzahl) ist diese Gruppe zyklisch von
  gerader Ordnung $p-1$ und besitzt daher genau ein Element der Ordnung $2$,
  nämlich $-1$.  Also $\prod_{g \in G} g = -1 \pmod{p}$. \qedhere
\end{proof}

% ================================================================
\section{Schlussfolgerung}
% ================================================================

Wir haben einen vollständigen elementaren Beweis des Satzes von Wilson und vier
natürliche Verallgemeinerungen vorgestellt.
Der elementare Charakter der Beweise verdeutlicht die tiefe Verbindung zwischen
der Struktur multiplikativer Gruppen modulo $n$ und der Arithmetik der Fakultäten.

Der Wilson-Quotient und Wilson-Primzahlen bleiben ein aktives Gebiet der
algorithmischen Zahlentheorie: Bekannt ist, dass keine Wilson-Primzahl
unterhalb von $2 \times 10^{13}$ existiert (Crandall--Dilcher--Pomerance, 1997),
und heuristische Argumente legen nahe, dass bis zu $N$ ungefähr $\ln\ln(N)$
Wilson-Primzahlen existieren.

\begin{thebibliography}{9}

\bibitem{Hardy1979}
G.~H.~Hardy und E.~M.~Wright.
\newblock \emph{Einführung in die Zahlentheorie} (An Introduction to the Theory of Numbers), 5.~Aufl.
\newblock Oxford University Press, 1979.

\bibitem{Ireland1990}
K.~Ireland und M.~Rosen.
\newblock \emph{A Classical Introduction to Modern Number Theory}, 2.~Aufl.
\newblock Springer, New York, 1990.

\bibitem{Crandall1997}
R.~Crandall, K.~Dilcher und C.~Pomerance.
\newblock A search for Wieferich and Wilson primes.
\newblock \emph{Math.\ Comp.}, 66(217):433--449, 1997.

\bibitem{Lagrange1771}
J.-L.~Lagrange.
\newblock Démonstration d'un théorème nouveau concernant les nombres premiers.
\newblock \emph{Nouv.\ Mém.\ Acad.\ Berlin}, 1771.

\end{thebibliography}

\end{document}
